\documentclass[A4paper,11pt]{article}
\usepackage{graphics,graphicx,epsfig,wrapfig}
\usepackage{longtable}  
\usepackage{amsmath,verbatim,amssymb,color,lscape}
\usepackage{kotex}
\usepackage{natbib}
\usepackage{lastpage}
\usepackage{multirow} 
\usepackage{authblk}
\usepackage{bm}
\usepackage{indentfirst}
\usepackage[ruled,vlined]{algorithm2e}
\usepackage{algorithmic}
\usepackage[normalem]{ulem}
\usepackage[export]{adjustbox}
%\usepackage{hyperref}
\usepackage{verbatim}
\usepackage{tabularx}
\usepackage{booktabs}
\usepackage{siunitx}

\sisetup{
  detect-mode,
  group-digits = false,
  table-number-alignment = center
}

\newcommand{\hide}[1]{}  
\newcommand{\msim}{\mathop{\rm \sim}}
\newcommand{\ind}{\msim\limits^{\mbox{\tiny ind}}}
\newcommand{\iid}{\msim\limits^{\mbox{\tiny iid}}}
\makeatletter
\newcommand{\s}{\vspace{.25cm}}
\renewcommand{\familydefault}{\sfdefault}
\usepackage[outerbars,color]{changebar}
\ifx\pdfoutput\undefined
\else\ifnum\pdfoutput>0
  \usepackage{pdfcolmk}
\fi\fi
\cbcolor{black}

\setlength{\changebarsep}{5mm}

\usepackage[margin=1.0in]{geometry}
\baselineskip=15.5pt

\usepackage{tikz}
\usetikzlibrary{bayesnet}
\usetikzlibrary{fit,positioning}
\newcommand{\obs}{latent}
\newcommand{\lat}{obs}
\newcommand{\con}{const}
\newcommand{\independent}{\perp\mkern-9.5mu\perp}

\setcounter{secnumdepth}{4} 
\renewcommand{\baselinestretch}{1.5}

\newfont{\rmm}{cmr10 at 11pt}
\rmm
\newcommand{\widesim}[2][1.5]{
  \mathrel{\overset{#2}{\scalebox{#1}[1]{$\sim$}}}
}

\pagestyle{plain}

\title{A Representation-Learning Item Response Model for Identifying Behaviorally Important Actions in PIAAC Process Data}

\author[1]{Junyeong Park}
\author[1]{Daeun Hwangbo}
\author[1]{Seyoung Park}
\author[1,2]{Ick Hoon Jin}
\author[3]{Minjeong Jeon}
\affil[1]{Department of Statistics and Data Science, Yonsei University. Republic of Korea.}
\affil[2]{Department of Applied Statistics, Yonsei University. Republic of Korea.}
\affil[3]{School of Education and Information Studies, University of California, Los Angeles. USA.}
\date{}

\begin{document}
\maketitle

\begin{abstract}
Problem-solving log process data collected in online environments contain rich information, including item difficulty, respondent ability, and problem-solving strategies. However, such log process data are complex and noisy, making it challenging to identify behaviors that substantially contribute to problem-solving. In this paper, we propose an representatoin-learning item response theory (IRT) model that incorporates a new behavioral term by encoding raw log data through hierarchical relationships with surrounding behaviors using an embedding framework from natural language processing. We then apply the spike-and-slab prior to selectively identify behaviors that significantly influence correct or incorrect responses. We apply the resulting framework to data from the OECD Programme for the International Assessment of Adult Competencies (PIAAC) for measuring problem solving in technology-rich environments (PSTRE).
\end{abstract}


\noindent {\bf Keywords:} Process data, Item Response Theory, Behavioral embedding, Spike-and-slab prior, Problem-solving assessment
\newpage


\section{Introduction}

Large-scale computer-based assessments increasingly record detailed behavioral logs of respondents' interactions with the test interface. These process data capture the sequence and timing of every action taken while solving a problem, from navigating the interface and inspecting available information to entering responses and making corrections \citep{greiff2016understanding, kroehne2018conceptualize, goldhammer2021byproduct, maddox2023uses, chen2024understanding}. The OECD Programme for the International Assessment of Adult Competencies (PIAAC) and, in particular, its problem solving in technology-rich environments (PSTRE) domain represents one of the richest sources of such data: each respondent-item pair yields an ordered sequence of discrete behavioral events together with exact timestamps that record the full problem-solving episode rather than merely its outcome \citep{oecd2012literacy, piaac2009pstre, organisation2019beyond, goldhammer2020analysing}.

The educational value of these process data has been widely recognized. Process data can reveal information about problem-solving strategies, use of available resources, and the trajectory of decision-making that is not captured by the binary or polytomous response outcome alone \citep{stadler2020first, eichmann2020exploring, he2021leveraging, reis2021improving, hahnel2023patterns}. For example, respondents who ultimately answer an item correctly may traverse systematically different sequences of interface interactions than those who answer it incorrectly, and the nature of these differences can inform both the interpretation of assessment results and the design of instructional feedback \citep{he2015identifying, he2016analyzing, tang2020latent, tang2021exploratory, xiao2021exploring, zhang2024external}.

Despite the interpretive richness of process data, most existing approaches to their analysis rely on manually constructed process indicators. Researchers typically select a small number of behavioral features --- such as the number of interface switches, the total time on task, the number of actions taken, or the occurrence of specific event types --- and use these as covariates in regression or latent variable models \citep{he2015identifying, he2016analyzing, goldhammer2020analysing}. This feature-engineering strategy has produced informative findings; however, it is fundamentally constrained by analysts' prior knowledge of which behaviors are likely to be relevant. Features that were not anticipated at the design stage may be omitted, and the selection of indicators introduces subjectivity that limits the reproducibility and generalizability of the analysis. Furthermore, summarizing a variable-length, ordered action sequence into a fixed set of scalar counts discards the sequential structure of the data, which may itself carry meaningful information about how respondents approach the task \citep{he2019clustering, tang2020latent, tang2021exploratory, wang2023subtask, ulitzsch2023machine, zhou2024investigating}.

A complementary strand of research has applied item response theory (IRT) to model response data from large-scale assessments \citep{rasch1960probabilistic, lord1980applications, vanderLinden2016handbook}. The Rasch model and its extensions provide a well-established framework for jointly estimating respondent ability and item difficulty from binary or polytomous responses. Recent work has explored ways to incorporate auxiliary information --- such as response times, eye-tracking data, and collateral test-taking behaviors --- into the IRT framework to enrich the measurement of latent constructs. Extending this approach to the full action sequences available in process data, however, requires a representation of behavioral sequences that is both compact enough for inclusion in a response model and sufficiently expressive to capture meaningful variation across respondents.

In this paper, we propose a three-stage framework that addresses these limitations by converting raw log sequences into latent action representations and incorporating them into an extended Rasch model. In the first stage, each action unit is embedded into a dense continuous vector using a hybrid Word2Vec model \citep{mikolov2013efficient, mikolov2013distributed} that adapts the skip-gram training objective to action sequences and incorporates a FastText-style composition \citep{bojanowski2017enriching} over the component tokens of each action label. The resulting embeddings capture the sequential co-occurrence structure and the internal semantic structure of the log event labels simultaneously. In the second stage, each respondent's augmented action sequence --- combining action embeddings with temporal features --- is reduced to a low-dimensional latent representation using an LSTM autoencoder \citep{hochreiter1997long, hinton2006reducing}. The per-action latent values preserve the sequential context of each behavioral event in a form that can be entered into the response model. In the third stage, the latent action representations are incorporated into the Rasch model as an additional term whose coefficient, the action weight, is assigned a spike-and-slab prior \citep{mitchellbeauchamp1988, george1993variable}. The spike-and-slab formulation enables simultaneous estimation and variable selection, yielding a posterior inclusion probability for each action that quantifies the evidence that it contributes to response accuracy beyond the contributions of respondent ability and item difficulty.

We apply the proposed framework to log process data from 14 PSTRE items in the United States PIAAC study. We demonstrate that the framework identifies important actions in a manner consistent with the task structure and reveals behavioral differences between correct and incorrect respondents that are not apparent from action frequencies alone. A simulation study confirms that the spike-and-slab representation-learning action-IRT model can recover truly important actions under data conditions that approximate the empirical application.

The main contributions of this paper are as follows. First, we propose an end-to-end framework for identifying important actions in problem-solving process data that does not require manual selection of process indicators, enabling systematic and reproducible behavioral analysis. Second, we introduce a hybrid Word2Vec model for action embedding that exploits both the sequential context of each action and the compositional structure of its label, providing a richer representation than standard bag-of-words or one-hot encodings. Third, we propose representation-learning action-IRT, an extension of the Rasch model that incorporates latent action representations as an additional predictor of response accuracy within the established psychometric IRT framework. Fourth, we apply the spike-and-slab prior within this model to perform automatic and data-driven variable selection over the action space, yielding interpretable posterior inclusion probabilities for each action-item combination.

The remainder of the paper is organized as follows. Section~2 reviews related work on process data analysis, sequential representation learning, and Bayesian IRT modeling. Section~3 describes the PIAAC data and preprocessing pipeline. Section~4 presents the proposed methodology. Section~5 presents the empirical results, including cross-item patterns and item-level case studies. Section~6 reports a simulation study evaluating the variable-selection performance of the representation-learning action-IRT model. Section~7 discusses the implications and limitations of the findings. Section~8 concludes.

\section{Background and Related Work}

\subsection{PIAAC PSTRE and Process Data}\label{sec:piaac_data}

The Programme for the International Assessment of Adult Competencies (PIAAC) is a large-scale international assessment developed by the Organisation for Economic Co-operation and Development (OECD) to measure proficiency in key information-processing skills among adults aged 16--65; its first cycle was conducted in 2011--2012 \citep{oecd2013piaac}. Among the three domains assessed---literacy, numeracy, and problem solving in technology-rich environments (PSTRE)---the PSTRE domain is of particular relevance to process data research because it requires respondents to interact with simulated digital environments, such as email clients, web browsers, and spreadsheets, to complete realistic problem-solving tasks \citep{piaac2009pstre}. A distinguishing feature of the computer-based PSTRE assessment is the systematic collection of process data through detailed log files. These logs consist of timestamped records of respondents' interactions with the testing application, capturing each discrete action --- such as page navigation, drag-and-drop operations, and form submissions --- as it occurs during task completion. Unlike traditional outcome measures that record only whether a response is correct, these behavioral records provide a fine-grained trace of how individuals approach and execute problem-solving tasks, encoding the trajectory of problem-solving behavior rather than merely its final outcome.

The raw log data are organized at the event level; each row corresponds to a single recorded event and contains fields including the country identifier, respondent sequence identifier (\texttt{SEQID}), booklet identifier, item identifier, event name, event type, timestamp, and event description. The field \texttt{event\_type} provides a broad action category (e.g., \texttt{MAIL\_DROP}, \texttt{FOLDER\_VIEWED}, \texttt{TOOLBAR}), while \texttt{event\_description} supplies contextual detail about the specific object or interface element involved---for example, identifying which email message was moved and to which destination folder. Because a single user action may generate multiple log entries, including both core actions initiated by the respondent and ancillary system-generated events, the raw logs require careful preprocessing to construct analysis-ready action sequences. The preprocessing pipeline applied in this study is described in Section~\ref{sec:preprocessing_pipeline}.

\begin{table}[htbp]
\centering
\caption{PSTRE items used in the analysis. $N$ denotes the number of respondent--item pairs with non-empty action sequences after preprocessing.}
\label{tab:items}
\begin{tabular}{llcr}
\toprule
Item & Task & Response type & $N$ \\
\midrule
PS1-1 & Party Invitations (Part~1) & Polytomous (0--3) & 1,295 \\
PS1-2 & Party Invitations (Part~2) & Binary            & 1,246 \\
PS1-3 & CD Tally                   & Binary            & 1,272 \\
PS1-4 & Sprained Ankle (Part~1)    & Binary            & 1,255 \\
PS1-5 & Sprained Ankle (Part~2)    & Binary            & 1,302 \\
PS1-6 & Tickets                    & Binary            & 1,282 \\
PS1-7 & Class Attendance           & Polytomous (0--3) & 1,074 \\
\midrule
PS2-1 & Club Membership (Part~1)   & Binary            & 1,274 \\
PS2-2 & Club Membership (Part~2)   & Polytomous (0--3) & 1,170 \\
PS2-3 & Book Order                 & Binary            & 1,238 \\
PS2-4 & Meeting Room               & Polytomous (0--3) & 1,161 \\
PS2-5 & Reply All                  & Binary            & 1,199 \\
PS2-6 & Locate Email               & Polytomous (0--3) & 1,131 \\
PS2-7 & Lamp Return                & Polytomous (0--3) & 1,230 \\
\bottomrule
\end{tabular}
\end{table}

The present analysis uses log process data from the U.S.\ sample of the first PIAAC cycle for the 14 PSTRE items. These 14 items are delivered in simulated digital environments and span a broad range of task complexity. Simpler items require single-application operations---such as sorting emails into designated folders or selecting a value from a spreadsheet---whereas more complex items demand navigation across multiple applications, evaluation of information from several sources, and integration of findings into a coherent solution. For instance, the Party Invitations item (U01a) involves a relatively straightforward email-management task, while the Lamp Return item (U23) requires respondents to navigate a multi-page website, complete an online exchange form, and confirm the transaction. Some items are scored as binary correct/incorrect outcomes, whereas others employ a polytomous scoring scheme on a 0--3 scale. Table~\ref{tab:items} summarizes the 14 items, their scoring formats, and the number of respondent--item pairs retained after preprocessing.

\subsection{Prior Approaches to Process Data Analysis}

Early work on PIAAC process data primarily adopted feature-engineering approaches, constructing scalar summary statistics from the log sequences and incorporating them into regression or latent variable models. Frequently used features include the number of actions, the total time on task, the number of interface switches between applications, and indicators for specific event types such as tool usage or response revision \citep{he2015identifying, he2016analyzing, goldhammer2020analysing}. These features have been shown to carry meaningful associations with response accuracy and to predict performance beyond the level explained by sociodemographic covariates. However, the feature set is determined by the analyst prior to modeling, which introduces subjectivity and restricts the analysis to behaviors anticipated at the design stage. More recent work has applied sequence mining, clustering, and machine learning methods to extract behavioral patterns more systematically from the raw event streams \citep{ulitzsch2023machine}. While these methods reduce reliance on pre-specified features, they often produce representations that are difficult to interpret within a psychometric framework or that do not readily yield action-level effect estimates that are comparable across items and respondents.

A complementary line of research has extended IRT-based models to incorporate process information alongside the response outcome. \citet{he2019clustering} proposed a joint model linking response accuracy, response time, and process indicators through a shared latent ability construct. \citet{tang2020latent} developed a latent class model that jointly describes response patterns and behavioral sequences. These approaches represent important steps toward integrating process data into psychometric modeling; however, they continue to depend on pre-specified behavioral summaries rather than learning action representations directly from the raw sequences.

Taken together, these strands of research reveal a recurring tension in the analysis of assessment process data: methods that remain interpretable within a psychometric framework tend to rely on behavioral summaries specified in advance by the analyst, whereas methods that learn behavioral patterns directly from the raw event streams tend to produce representations that are difficult to relate to response accuracy at the level of individual actions. The framework developed in this paper is positioned between these two traditions. It learns action representations directly from the raw log sequences, dispensing with pre-specified process indicators, and embeds these representations within an extended Rasch model so that the contribution of each action to response accuracy can be estimated and interpreted within the established IRT framework. The two methodological foundations on which it builds---representation learning for sequential data and Bayesian item response modeling with variable selection---are reviewed in the remainder of this section.

\subsection{Representation Learning for Sequential Data}

Word2Vec introduced the idea of learning dense vector representations of discrete symbols from their distributional context in large corpora \citep{mikolov2013efficient, mikolov2013distributed}. The skip-gram model trains a shallow neural network to predict the surrounding symbols within a fixed context window given a center symbol, producing embeddings in which symbols that appear in similar contexts have similar vector representations. FastText extended the skip-gram framework by decomposing each word into character $n$-grams and representing the word as the sum of its unit-level and subword-level embeddings \citep{bojanowski2017enriching}. This subword composition allows the model to generalize across morphologically related words and to produce useful representations for rare or unseen words, which is especially valuable when the vocabulary is large and individual items are infrequent.

Embedding methods originally developed for natural language processing have been successfully adapted to a wide range of sequential data domains. In electronic health records, medical codes have been embedded from co-occurrence patterns in patient visit sequences to support downstream clinical prediction tasks \citep{choi2016multi}. In user behavior modeling, click sequences, purchase histories, and navigation logs have been embedded using skip-gram-style objectives to learn distributed representations of user actions and items \citep{grbovic2015commerce, barkan2016item2vec}. In educational data mining, knowledge concepts, learning activities, and student interaction sequences have been represented as dense embeddings to improve models of knowledge acquisition and student performance \citep{pardos2020university}. These applications demonstrate that the skip-gram framework and its extensions provide a flexible and effective strategy for encoding the co-occurrence structure of discrete sequential data, independent of the natural-language domain in which the method was originally developed.

Long short-term memory (LSTM) networks are a class of recurrent neural networks designed to model long-range dependencies in sequential data through gated memory cells \citep{hochreiter1997long}. An LSTM encoder processes an input sequence step by step, updating a hidden state that summarizes the history of the sequence up to each time point. The autoencoder architecture pairs an LSTM encoder with a corresponding decoder that reconstructs the input sequence from the encoded representation; the bottleneck layer at the junction of encoder and decoder learns a compressed latent representation that captures the essential structure of the sequence \citep{srivastava2015unsupervised}. LSTM autoencoders have been applied to time series, activity recognition, and anomaly detection tasks, where they learn compact representations of variable-length sequential observations without relying on predefined feature extraction \citep{malhotra2016lstm}. In the proposed framework, the LSTM autoencoder serves the specific purpose of compressing each augmented action embedding vector into a low-dimensional latent value that can be incorporated as a predictor in the item response model.

\subsection{Bayesian Item Response Models with Auxiliary Information}

Item response theory provides a family of probabilistic models for relating a latent ability trait to the probability of a correct response on a set of test items \citep{rasch1960probabilistic, lord1980applications, vanderLinden2016handbook}. The Rasch model, the simplest member of this family, parameterizes the log-odds of a correct response as the difference between respondent ability and item difficulty. Bayesian inference for IRT models, typically implemented via Markov chain Monte Carlo, allows the posterior distribution of all model parameters to be estimated simultaneously, supports the incorporation of prior knowledge through the prior distribution, and facilitates model comparison through posterior predictive checks \citep{patz1999straightforward}.

Several extensions of the Rasch model have been proposed to incorporate information beyond the response outcome. Response time models link the latent speed of a respondent to the time required to complete an item, typically through a log-normal model for response times and a Rasch model for response accuracy, with a hierarchical structure connecting speed and ability \citep{van2007hierarchical}. Mixture IRT models allow the item parameters or the relationship between ability and response to differ across latent subgroups, accommodating behavioral heterogeneity not captured by a single latent trait \citep{rost1990rasch}. More recently, approaches that incorporate collateral information---such as demographic variables, response strategies, or process indicators---have been used to augment the standard IRT likelihood, improving the precision of ability estimates or providing estimates of the marginal effect of specific behavioral features \citep{wang2015mixture}.

Bayesian variable selection provides a principled framework for determining which predictors in a regression model have non-negligible associations with the response variable. The spike-and-slab prior, introduced by \citet{mitchellbeauchamp1988} and extended by \citet{george1993variable}, assigns each regression coefficient a two-component mixture: a point mass at zero (the spike) representing the hypothesis that the coefficient is effectively zero, and a diffuse distribution (the slab) representing the hypothesis that the coefficient has a substantive effect. The posterior inclusion probability for each coefficient---the posterior probability that it belongs to the slab component---provides a natural summary of the evidence for its inclusion in the model. This framework has been applied to high-dimensional variable selection problems in statistics \citep{ishwaran2005spike} and has been used in psychometric models to identify items or covariates that contribute to latent variable models \citep{chen2022advantages}. In the proposed representation-learning action-IRT model, the spike-and-slab prior is placed on the action weights to enable automatic, data-driven selection of the actions that are associated with the probability of a correct response, without requiring the analyst to pre-specify which actions are expected to be relevant.


%=============Preprocessing==================================================================================
\section{Action Unit Construction and Preprocessing}
\subsection{Unit and Token Definitions}\label{sec:unit_token}

The central object in this study is the \emph{action unit}. An action unit is a standardized behavioral label that represents one meaningful event in a problem-solving sequence of a respondent. Each action unit is constructed by combining the broad event category recorded in \texttt{event\_type} with, when available, a cleaned portion of the contextual information recorded in \texttt{event\_description}. The resulting label preserves both the type of interface operation performed and the task-specific detail needed to distinguish different instances of the same operation.

Each action unit has an internal compositional structure. A unit label consists of semantic components connected by underscores, and these components are referred to as \emph{action tokens}. In the example \texttt{mail-drop\_cancome-folder}, the tokens \texttt{mail-drop}, and \texttt{cancome-folder} correspond respectively to the application context, the operation type, and the target object. This decomposition is not merely notational: it reflects the hierarchical organization of the log event labels, in which a higher-level operation category is progressively qualified by lower-level descriptors that specify the object or location of the interaction.

The distinction between units and tokens is important for the action embedding model described in Section~\ref{sec:embedding}. The unit serves as the behavioral element in the action sequence and defines the vocabulary over which sequential co-occurrence patterns are learned. The tokens provide sub-unit semantic components that enable information sharing across related units. Two units that differ in their target object but share tokens indicating the same interface operation---for example, \texttt{mail-drop\_cancome-folder} and \texttt{mail-drop\_cannotcome-folder}---receive partially shared representations through their common token components (\texttt{mail-drop}), even if one unit is rare and the other is frequent. This compositional structure motivates the FastText-style subword composition adopted in Section~\ref{sec:embedding}.

\subsection{Preprocessing Pipeline}\label{sec:preprocessing_pipeline}

The raw log files contain events that are too granular, redundant, or system-dependent to serve directly as modeling units. We therefore preprocess the logs to transform raw event records into consistent action sequences while preserving the observed temporal order and meaningful behavioral distinctions. The preprocessing procedure follows the guidance provided in \citet{hwangbo2026tutorial}, which describes common preprocessing issues in PIAAC process data and practical strategies for addressing them. We adopt the three-stage structure recommended therein---rule-based cleaning, LLM-assisted description standardization, and action unit construction---and adapt each stage to the analytic aims of the present study. The complete preprocessing scripts are available at (Github).

The first stage consists of rule-based preprocessing performed by the researcher. This stage addresses three categories of issues identified in \citet{hwangbo2026tutorial}: timestamp reversals, duplicate records, and multiple log entries per user action. Restart events require special handling because the system-generated restart resets the timestamp to zero in the middle of a sequence. Following the offset-correction procedure described in \citet{hwangbo2026tutorial}, timestamps recorded after a restart are adjusted by adding a restart-specific offset so that the full sequence remains temporally ordered on a single cumulative time scale. Duplicate records---multiple log entries with the same event type and timestamp within the same respondent--item pair---are removed. Events that are administrative or system-generated rather than respondent-initiated are also excluded; examples include the session-start marker recorded at item presentation and navigation confirmation clicks that do not represent substantive problem-solving actions. In addition, when a single user action generates multiple log entries comprising a core action and accompanying ancillary events, these entries are consolidated into a single representative action that retains the analytically relevant information. Finally, consecutive \texttt{KEYPRESS} events are collapsed into a single keypress action whose label records the number of consecutive keypresses, preventing long runs of individual keystrokes from dominating the action sequence while preserving information about the extent of text entry.

As part of this stage, action labels that recur across items with inconsistent or uninformative raw identifiers are standardized. For example, submission-related interface elements are described differently across items in the raw data; these are mapped to consistent labels such as "next", "OK", and "cancel" to ensure comparability in the downstream embedding model.

The second stage addresses the event descriptions that remain after rule-based cleaning. Although the first stage resolves structural issues in the log files, the \texttt{event\_description} field often retains system-specific identifiers, numerical indices, and delimiter-separated strings that are too noisy or inconsistent to use directly as components of action unit labels. These descriptions must therefore be normalized into concise substitutes that preserve the core distinguishing information---for instance, by reducing a raw description containing internal element identifiers to a label that retains only the target object or interface location. Because this normalization task involves context-dependent string interpretation across a large and heterogeneous set of descriptions, we employ an LLM-assisted procedure adapted from the prompt-based workflow introduced in \citet{hwangbo2026tutorial}. The procedure is applied within event-type groups so that descriptions from different action categories are processed independently. The design of the prompt, the choice of model, and the specific adaptations made for the present study are described in Section~\ref{sec:llm_preprocessing}.

The third stage combines the cleaned event type and the cleaned description into the final action unit, following the construction rule defined in Section~\ref{sec:unit_token}. If the cleaned description contains meaningful information, it is concatenated with the event type using underscores. If the description is empty or not substantively informative, the event type alone is retained as the unit label. The resulting data for each item consist of respondent-level ordered sequences of action units and their associated timestamps. These sequences serve as the input to the action embedding model described in Section~\ref{sec:embedding}.

\subsection{LLM-Assisted Description Standardization}\label{sec:llm_preprocessing}

The LLM-assisted step normalizes the detailed event descriptions that remain after rule-based cleaning. Its scope is restricted to string-level normalization rather than inference about respondents' latent strategies or substantive analytic decisions.

The input to the LLM is organized as groups of descriptions within each \texttt{event\_type}. For each group, the model receives the set of raw description strings and returns a structured table containing the original event type, the original description, and a proposed concise substitute that retains the core distinguishing information. Several constraints are imposed to keep the transformation conservative: task-specific identifiers are preserved as single semantic units, redundant components are removed only when they duplicate information already present elsewhere in the label, and descriptions that do not form a coherent group are retained in their original form. The complete set of transformation rules and constraint specifications is provided together with the full prompt text in Section~B of the Supplementary Material.

After the LLM-generated substitutes are produced, the researcher reviews each mapping and connects the cleaned substitute values to the corresponding event types to form final action units. This human-in-the-loop design uses the LLM for scalable normalization of a large and heterogeneous set of description strings while ensuring that the final preprocessing decisions remain auditable and under researcher control. In the present analysis, Claude Sonnet 4.5 \citep{anthropic2025sonnet} is used for this step, accessed via the OpenRouter API.

\section{Methodology}

% TODO: Insert a schematic figure (e.g., Figure~\ref{fig:framework}) illustrating the three-stage
%       pipeline: (1) action embedding, (2) LSTM autoencoder, (3) representation-learning action-IRT with spike-and-slab.
%       The Table of Contents specifies a flowchart for this subsection.

This section describes the proposed modeling framework for identifying discriminative actions in problem-solving process data. A discriminative action is one whose latent representation differs systematically between correct and incorrect respondents after controlling for respondent ability and item difficulty. The framework consists of three stages applied after preprocessing. First, each action unit is mapped to a continuous vector through a hybrid Word2Vec embedding model that uses both action-unit sequences and their component tokens. Second, the action embedding sequence of each respondent is augmented with temporal features and compressed into a low-dimensional latent representation using an LSTM autoencoder. Third, the latent action representations are incorporated into an extended Rasch model with a spike-and-slab prior that performs data-driven selection of discriminative actions among them.

\begin{figure}[htbp]
    \centering
    \includegraphics[width=0.95\linewidth]{figure/main_methodology.png}
    \caption{Main Framework}
    \label{fig:placeholder}
\end{figure}

Let $i=1,\ldots,N$ index respondents and $j=1,\ldots,J$ index items. For respondent $i$ and item $j$, let $Y_{ij}\in\{0,1\}$ denote the binary response outcome, where $Y_{ij}=1$ indicates a correct response. The preprocessed action sequence for the pair $(i,j)$ is
\[
\mathcal{A}_{ij}=(a_{ij1},a_{ij2},\ldots,a_{ijN_{ij}}),
\]
where $N_{ij}$ is the length of the sequence and each $a_{ijn}$ is an action unit. An action unit can be decomposed into lower-level action tokens that reflect the hierarchical structure of the log event label (see Section~3.2). For item $j$, let $N_{.j.}$ denote the number of unique action units observed across all respondents, and let $A_{ij}\subseteq\{1,\ldots,N_{.j.}\}$ be the set of unique action indices that appear in $\mathcal{A}_{ij}$.

The framework produces a latent action representation $C_{ijn}^{(d)}\in\mathbb{R}$ for each action occurrence $n$ and latent dimension $d=1,\ldots,D$. When a unique action $l$ occurs $R_{ijl}$ times in the sequence of respondent $i$ for item $j$, we summarize its representation by the within-sequence average
\[
\bar{C}_{ijl}^{(d)}
=\frac{1}{R_{ijl}}\sum_{k=1}^{R_{ijl}} C_{ijlk}^{(d)},
\quad d=1,\ldots,D,
\]
where $C_{ijlk}^{(d)}$ is the $d$th latent value of the $k$th occurrence of action $l$. 

\subsection{Hybrid Word2Vec for Action Embedding}\label{sec:embedding}

The first stage converts each action unit into a dense continuous vector. The key modeling assumption underlying this stage is that consecutive actions in a problem-solving sequence are stochastically dependent. Specifically, we assume that the action sequence $\mathcal{A}_{ij}$ exhibits approximate Markov dependence: conditional on the current action, the distribution of nearby actions is largely determined by a local context window rather than by the full history of the sequence. This assumption is plausible for PIAAC log data, in which an action at a given point in the problem-solving process---such as opening a folder---constrains the actions that can immediately follow, such as viewing or dragging a message within that folder. To encode this local sequential dependence, we adapt the skip-gram objective from Word2Vec \citep{mikolov2013efficient, mikolov2013distributed} to action sequences: a center action unit is used to predict neighboring action units within a fixed context window so that actions appearing in similar local contexts receive similar vector representations.

The PIAAC log data contain action labels with an internal hierarchical structure. As described in Section~3.2, each preprocessed action unit can be decomposed into lower-level action tokens that correspond to semantic components of the event label. For example, the unit \texttt{mail-drop\_cancome-folder} decomposes into tokens related to the mail operation, the drop action, and the target folder. Different action units may share one or more tokens while differing in others and this partial overlap carries semantic information that a purely unit-level embedding would discard.

To preserve this compositional structure, we adopt a FastText-style subword composition strategy \citep{bojanowski2017enriching}. The model maintains separate embedding vectors for both action units and action tokens. For an action unit $u$ with component tokens $g_1,\ldots,g_m$, the composite representation is
\[
\mathbf{e}(u)
= \mathbf{v}_{u}+\sum_{r=1}^{m}\mathbf{v}_{g_r},
\]
where $\mathbf{v}_{u}\in\mathbb{R}^{D_{\mathrm{Action}}}$ is the unit-level embedding and $\mathbf{v}_{g_r}\in\mathbb{R}^{D_{\mathrm{Action}}}$ is the embedding of token $g_r$. This additive composition combines the holistic distributional information of the action unit with compositional information shared across units that contain common tokens. In particular, rare action units that share tokens with more frequent units benefit from the token-level training signal, mitigating the data sparsity that would affect a purely unit-level model.

Training pairs are generated in two ways. The first type consists of \emph{unit--unit} pairs, in which the center unit is paired with each neighboring unit within the context window. The second type consists of \emph{token--unit} pairs, in which each token of the center unit is individually paired with the same neighboring units. Both pair types are trained jointly under a skip-gram objective with negative sampling. This dual training scheme encourages action units that occur in similar sequential contexts, or that share informative token components, to have similar vector representations.

After training, each action occurrence is represented by a vector in $\mathbb{R}^{D_{\mathrm{Action}}}$. For respondent $i$ and item $j$, the full action sequence yields an action embedding matrix
\[
\mathbf{E}_{ij}\in\mathbb{R}^{N_{ij}\times D_{\mathrm{Action}}}.
\]

The action embedding captures the sequential and compositional structure of the log events but does not encode when each action occurred. Because the temporal profile of problem-solving behavior carries information beyond the sequential order of actions, we augment each embedding vector with three temporal features. Let $t_{ijn}$ denote the
elapsed time from the start of the problem-solving episode to the $n$th action. The augmented vector is
\[
\tilde{\mathbf{e}}_{ijn}
=
\bigl[
\mathbf{e}(a_{ijn}),\;
t_{ijn},\;
t_{ijn}^{2},\;
t_{ijn}-t_{ij,n-1}
\bigr].
\]
Each temporal component serves a distinct purpose. The raw elapsed time $t_{ijn}$ locates the action within the overall timeline of the episode, allowing the downstream model to distinguish early-stage exploration from late-stage confirmation. Its squared term $t_{ijn}^{2}$ captures potential nonlinear associations between absolute timing and
behavioral quality---for example, actions performed very early or very late in an episode may carry different diagnostic information than those performed at intermediate times. The inter-action interval $t_{ijn}-t_{ij,n-1}$ encodes the local pacing of behavior: short intervals may indicate rapid, routine operations, whereas long intervals
may reflect deliberation, confusion, or off-task behavior. For the first action in each sequence ($n=1$), the inter-action interval is set to zero.

Because the raw temporal variables operate on a time scale (seconds) that can differ by orders of magnitude from the embedding coordinates, directly concatenating unscaled timestamps could allow the temporal features either to dominate the input to the downstream LSTM autoencoder or, conversely, to become numerically negligible. To ensure that the
temporal features contribute on a comparable numerical scale, each of the three temporal variables is standardized within item $j$ to have zero mean and unit variance across all action occurrences observed for that item. Let $f:\mathbb{R}\to\mathbb{R}$ denote item-level standardization. The augmented embedding vector for the $n$th action is
\[
\tilde{\mathbf{e}}_{ijn}
=
\bigl[
\mathbf{e}(a_{ijn}),\;
f(t_{ijn}),\;
f(t_{ijn}^{2}),\;
f(t_{ijn}-t_{ij,n-1})
\bigr].
\]
The resulting augmented sequence matrix is
\[
\tilde{\mathbf{E}}_{ij}\in\mathbb{R}^{N_{ij}\times(D_{\mathrm{Action}}+3)}.
\]

\subsection{LSTM AutoEncoder for Dimension Reduction}

The second stage reduces the augmented action embedding vectors to low-dimensional latent values that serve as action-level covariates in the item response model. This reduction is motivated by two considerations. First, the augmented embedding $\tilde{\mathbf{e}}_{ijn}\in\mathbb{R}^{D_{\mathrm{Action}}+3}$ is too high-dimensional to enter the response model directly. Assigning a separate regression weight to each embedding coordinate of every unique action would introduce a number of parameters far exceeding the available sample size, rendering estimation unstable and interpretation impractical. A low-dimensional summary of each action occurrence is therefore required. Second, the skip-gram embedding and its temporal augmentation encode the identity, compositional structure, and timing of each action, but they do not capture the \emph{sequential context} in which the action occurs. Under the approximate Markov assumption adopted in Section~4.2, the embedding reflects local co-occurrence patterns, yet problem-solving behavior often exhibits longer-range dependencies: the same action may carry different behavioral significance depending on whether it appears early in an exploratory phase or late in a confirmatory phase of the solution process. A dimensionality-reduction method that accounts for the ordering of the full action sequence can encode this contextual information into the latent representation of each action. These two requirements---dimensionality reduction and sequential context encoding---motivate the use of a Long Short-Term Memory (LSTM) autoencoder \citep{hochreiter1997long, srivastava2015unsupervised}. The LSTM architecture maintains a hidden state that is updated recurrently as each action in the sequence is processed, allowing the encoder to summarize the history of preceding actions at every time step. The autoencoder training objective ensures that the resulting latent representation retains sufficient information to reconstruct the original input sequence, providing an unsupervised criterion for learning compact representations without requiring response labels.

For each item $j$, the input to the LSTM autoencoder is the collection of augmented embedding matrices $\{\tilde{\mathbf{E}}_{ij}\}_{i=1}^{N}$. Let $H_{\mathrm{enc}}$ denote the hidden dimension of the encoder. The encoder LSTM processes the sequence $(\tilde{\mathbf{e}}_{ij1},\ldots,\tilde{\mathbf{e}}_{ijN_{ij}})$ step by step and produces a hidden state $\mathbf{h}_{ijn}\in\mathbb{R}^{H_{\mathrm{enc}}}$ at each position $n$. A linear projection then maps each hidden state to the $D$-dimensional latent space:
\[
\mathbf{C}_{ijn}
=
\mathbf{W}_{\mathrm{proj}}\,\mathbf{h}_{ijn}+\mathbf{b}_{\mathrm{proj}}
\in\mathbb{R}^{D},
\]
where $\mathbf{W}_{\mathrm{proj}}\in\mathbb{R}^{D\times H_{\mathrm{enc}}}$ and $\mathbf{b}_{\mathrm{proj}}\in\mathbb{R}^{D}$ are learnable parameters. Because the hidden state $\mathbf{h}_{ijn}$ depends on all preceding inputs $\tilde{\mathbf{e}}_{ij1},\ldots,\tilde{\mathbf{e}}_{ijn}$, the latent value $\mathbf{C}_{ijn}$ encodes not only the $n$th action itself but also the sequential context in which it was performed. This property distinguishes the LSTM-based reduction from context-free alternatives, such as principal component analysis or standard autoencoders applied independently to each action vector. The decoder reverses the encoding process: it maps each latent vector back to the hidden dimension through a linear layer, processes the resulting sequence with a decoder LSTM, and reconstructs the original augmented embedding through a final linear projection. The model is trained by minimizing the mean squared reconstruction error over observed, non-padded positions:
\[
\mathcal{L}_{j}
=
\frac{1}{\sum_{i}N_{ij}}
\sum_{i=1}^{N}\sum_{n=1}^{N_{ij}}
\bigl\|\tilde{\mathbf{e}}_{ijn}
-\hat{\mathbf{e}}_{ijn}\bigr\|^{2},
\]
where $\hat{\mathbf{e}}_{ijn}$ is the reconstructed vector at position $n$. The summation runs only over observed sequence positions; padded positions introduced to accommodate variable sequence lengths within training batches are masked and excluded from the loss. Training is performed separately for each item $j$, ensuring that the latent representation is adapted to the behavioral vocabulary and sequence characteristics of each task. The resulting output for respondent $i$ and item $j$ is a latent matrix $\mathbf{C}_{ij}=(\mathbf{C}_{ij1},\ldots,\mathbf{C}_{ijN_{ij}})^{\top} \in\mathbb{R}^{N_{ij}\times D}$, which preserves the action-level sequence length while replacing each $(D_{\mathrm{Action}}+3)$-dimensional augmented embedding with a $D$-dimensional context-aware representation. The columns of $\mathbf{C}_{ij}$ correspond to the $D$ latent dimensions, and each row provides a latent summary of a single action occurrence that is used as input to the representation-learning action-IRT model in Section~4.4.

The latent values $\mathbf{C}_{ij}$ are treated as fixed, observed covariates in the subsequent item response model. This two-stage plug-in strategy is computationally convenient and avoids the need to specify a joint likelihood over both the reconstruction and the response outcomes, but it does not propagate the uncertainty of the dimension-reduction stage into the posterior inference of the representation-learning action-IRT parameters. To the extent that the LSTM autoencoder provides a stable and accurate summary of the input sequences---as assessed by reconstruction error on held-out data---the practical impact of this approximation is expected to be modest. A fully joint estimation framework that integrates the encoding and response-modeling stages represents a direction for future methodological development.

\subsection{Representation-learning Action-IRT Model with Spike-and-Slab Selection}

The final stage links the latent action representations to response accuracy within a psychometric framework. We extend the Rasch model by introducing an additive action term that depends on the latent representations of the actions observed during each respondent's problem-solving episode. This extension preserves the standard IRT decomposition into respondent ability and item difficulty while providing a mechanism to quantify the association between specific behavioral events and the probability of a correct response.

\subsubsection{Model Specification}

Let $\pi_{ij}=P(Y_{ij}=1)$ denote the probability that respondent $i$ answers item $j$ correctly. The proposed representation-learning action-IRT model specifies the log-odds of a correct response as
\begin{equation}\label{eq:action_irt}
\mathrm{logit}(\pi_{ij})
=
\alpha_i+\beta_j
+
\sum_{l=1}^{N_{.j.}}
\left(
\sum_{d=1}^{D}
\omega_{jl}^{(d)}\,\bar{C}_{ijl}^{(d)}
\right)
I(l\in A_{ij}),
\end{equation}
where $\alpha_i$ is the ability parameter for respondent $i$, $\beta_j$ is the difficulty parameter for item $j$, and $\omega_{jl}^{(d)}$ is the action weight associated with unique action $l$ in item $j$ along latent dimension $d$. The indicator function $I(l\in A_{ij})$ ensures that action $l$ contributes to the linear predictor only when it appears in the observed sequence $\mathcal{A}_{ij}$; actions not performed by respondent $i$ on item $j$ are excluded from the summation. The inner sum aggregates the contributions across latent dimensions, so that the total effect of action $l$ on the log-odds is $\sum_{d=1}^{D}\omega_{jl}^{(d)}\,\bar{C}_{ijl}^{(d)}$ when $l\in A_{ij}$ and zero otherwise.

The action weight $\omega_{jl}^{(d)}$ has a direct interpretation: it measures the change in the log-odds of a correct response associated with a one-unit increase in the average latent representation $\bar{C}_{ijl}^{(d)}$ of action $l$ along dimension $d$, conditional on respondent ability and item difficulty. Because the latent representations $\bar{C}_{ijl}^{(d)}$ are derived from an unsupervised encoding of the action sequence (Section~4.3), the sign and magnitude of $\omega_{jl}^{(d)}$ reflect data-driven associations rather than theoretically pre-specified effect directions.

\subsubsection{Prior Specification and Variable Selection}

Because the number of unique actions per item is large and most actions are expected to have negligible associations with response accuracy after conditioning on ability and difficulty, the model requires a mechanism for automatic variable selection over a potentially high-dimensional action space. To this end, we assign each action weight a spike-and-slab prior \citep{mitchellbeauchamp1988, george1993variable}:
\begin{equation}\label{eq:spike_slab}
\omega_{jl}^{(d)}\mid \lambda_{jl}^{(d)}
\sim
\begin{cases}
N(0,\,\tau^2), & \text{if } \lambda_{jl}^{(d)}=0 \quad (\text{spike}),\\[4pt]
N(0,\,\nu^2), & \text{if } \lambda_{jl}^{(d)}=1 \quad (\text{slab}),
\end{cases}
\end{equation}
where $\tau^2$ and $\nu^2$ are the spike and slab variances with $\tau^2 \ll \nu^2$. The spike component concentrates the prior mass of $\omega_{jl}^{(d)}$ near zero, shrinking the weights of non-discriminative actions toward negligibility, whereas the slab component is diffuse and permits substantively large weights for actions whose latent representations are associated with response accuracy. The latent binary indicator $\lambda_{jl}^{(d)}$ governs component membership and is assigned an independent Bernoulli prior,
\begin{equation}\label{eq:lambda_prior}
\lambda_{jl}^{(d)}\sim \mathrm{Bernoulli}(0.5),
\end{equation}
reflecting an uninformative prior belief that, before observing the data, each action is equally likely to be discriminative or non-discriminative. The posterior distribution of $\lambda_{jl}^{(d)}$ thus drives the data-driven selection of discriminative actions. To complete the model specification, the priors on the remaining model parameters are
\begin{equation}\label{eq:remaining_priors}
\alpha_i \mid \sigma_{\alpha}^{2} \sim N\!\left(0,\,\sigma_{\alpha}^{2}\right),
\qquad
\beta_j \sim N\!\left(0,\,\sigma_{\beta}^{2}\right),
\qquad
\sigma_{\alpha}^{2} \sim \text{Inverse-Gamma}\!\left(a_0,\,b_0\right),
\end{equation}
where the respondent ability $\alpha_i$ is assigned a normal prior with variance $\sigma_{\alpha}^{2}$, the item difficulty $\beta_j$ a normal prior with fixed variance $\sigma_{\beta}^{2}$, and the ability variance $\sigma_{\alpha}^{2}$ a conjugate inverse-gamma hyperprior. The specific values of the hyperparameters $\sigma_{\beta}^{2}$, $a_0$, and $b_0$, together with the spike and slab variances $\tau^2$ and $\nu^2$, are reported in Section C of the Supplementary Material.

\subsubsection{Posterior Inference}

Let $\boldsymbol{\theta}=\{\{\alpha_i\},\{\beta_j\},\{\omega_{jl}^{(d)}\},
\{\lambda_{jl}^{(d)}\},\sigma_{\alpha}^{2}\}$ denote the full collection of model
parameters. Conditional on the observed responses $\mathbf{Y}=\{Y_{ij}\}$ and the
fixed latent action representations $\{\bar{C}_{ijl}^{(d)}\}$ obtained from the LSTM
autoencoder, the likelihood of the representation-learning action-IRT model is
\begin{equation}\label{eq:likelihood}
L(\boldsymbol{\theta}\mid\mathbf{Y})
=
\prod_{(i,j)\in\mathcal{O}}
\pi_{ij}^{\,Y_{ij}}
(1-\pi_{ij})^{1-Y_{ij}},
\end{equation}
where $\mathcal{O}=\{(i,j):Y_{ij}\text{ is observed}\}$ is the set of observed
respondent-item pairs and $\mathrm{logit}(\pi_{ij})$ is given by~\eqref{eq:action_irt}.
Combining the likelihood with the prior distributions
in~\eqref{eq:spike_slab}--\eqref{eq:remaining_priors}, the joint posterior is
\begin{equation}\label{eq:joint_posterior}
p(\boldsymbol{\theta}\mid\mathbf{Y})
\;\propto\;
L(\boldsymbol{\theta}\mid\mathbf{Y})
\
\prod_{i=1}^{N}p(\alpha_i\mid\sigma_{\alpha}^{2})
\
\prod_{j=1}^{J}p(\beta_j)
\
\prod_{j,l,d}
p(\omega_{jl}^{(d)}\mid\lambda_{jl}^{(d)})
\,p(\lambda_{jl}^{(d)})
\
p(\sigma_{\alpha}^{2}).
\end{equation}

Direct sampling from~\eqref{eq:joint_posterior} is intractable, so we employ a
component-wise Markov chain Monte Carlo (MCMC) sampler in which each parameter block is
updated from its full conditional distribution. The ability, difficulty, and
action-weight parameters ($\alpha_i$, $\beta_j$, $\omega_{jl}^{(d)}$) enter the model
through the logistic link and have no closed-form full conditionals; these blocks are
updated using Metropolis--Hastings steps with Gaussian random-walk proposals. The
inclusion indicators $\lambda_{jl}^{(d)}$ and the ability variance $\sigma_{\alpha}^{2}$
have closed-form full conditionals---Bernoulli and inverse-gamma, respectively---and are
updated by Gibbs sampling. The explicit full conditional distributions, the update
scheme, the proposal specifications, and the convergence assessment are provided in
Section~C of the Supplementary Material.

\subsubsection{Identification of Discriminative Actions}

The posterior samples of $\lambda_{jl}^{(d)}$ yield the posterior inclusion probability (PIP) for each action-dimension combination, defined as $\mathrm{PIP}_{jl}^{(d)}=P(\lambda_{jl}^{(d)}=1\mid\text{data})$, which is estimated by the proportion of post-burn-in MCMC iterations in which $\lambda_{jl}^{(d)}=1$. A high PIP indicates that the data favor the slab component for that action weight, providing evidence that the action's latent representation is associated with response accuracy beyond the contributions of ability and difficulty. We identify the set of active latent dimensions for action $l$ in item $j$ as
\begin{equation}\label{eq:active_set}
S_{jl}=\bigl\{d\in\{1,\ldots,D\}:
\mathrm{PIP}_{jl}^{(d)}\geq 0.5\bigr\}.
\end{equation}
The PIP and the active set $S_{jl}$ constitute the variable-selection layer of the model: they identify which action weights are supported by the slab component. To translate a selected weight into a behaviorally interpretable quantity, we next contrast its contribution between respondents who answered the item correctly and those who did not.

Let $\mathcal{C}_{j}=\{i:Y_{ij}=1\}$ and $\mathcal{I}_{j}=\{i:Y_{ij}=0\}$ denote the sets of respondents who answered item $j$ correctly and incorrectly, respectively. For each group $G\in\{\mathcal{C}_{j},\mathcal{I}_{j}\}$, the group-level average latent representation of action $l$ along dimension $d$ is
\begin{equation}\label{eq:group_mean}
\bar{C}_{jl}^{G(d)}
=
\frac{1}{|G|}\sum_{i\in G}\bar{C}_{ijl}^{(d)}.
\end{equation}
At each post-burn-in MCMC iteration $t$, the group-specific action effect is computed by activating only the slab-selected dimensions within the active set:
\begin{equation}\label{eq:group_effect}
E_{jl}^{G(t)}
=
\sum_{d\in S_{jl}}
\omega_{jl}^{(d)(t)}\,
I\!\left(\lambda_{jl}^{(d)(t)}=1\right)\,
\bar{C}_{jl}^{G(d)}.
\end{equation}
This quantity combines the posterior draw of the action weight with the observed group-level latent representation, restricting the summation to dimensions that have been selected as active. The difference between the correct-group and incorrect-group effects,
\begin{equation}\label{eq:delta_E}
\Delta E_{jl}^{(t)}
=
E_{jl}^{\mathcal{C}_{j}(t)}
-
E_{jl}^{\mathcal{I}_{j}(t)},
\end{equation}
summarizes whether action $l$ contributes differently to the linear predictor for correct versus incorrect respondents on item $j$ at iteration $t$. The full posterior distribution of $\Delta E_{jl}$ is obtained by computing~\eqref{eq:delta_E} at every saved iteration. Whereas the PIP identifies actions whose weights are supported by the data, $\Delta E_{jl}$ provides the criterion for the final, interpretation-oriented classification: we classify action $l$ as a \emph{discriminative action} for item $j$ when the 95\% highest posterior density (HPD) interval of $\Delta E_{jl}$ does not contain zero. The sign, magnitude, and posterior uncertainty of $\Delta E_{jl}$ together provide the basis for the substantive interpretation of each discriminative action reported in Section~6.

%%%%%%%%%%%%%%%%%%%%%%%%%%%%%%%%%%%%%%%%%%%%%%%%
\section{Real Data Application Results}

The empirical analysis used the one-dimensional latent representation ($D=1$) with robust scaling of the reduced action values applied to the data. The posterior was sampled using the component-wise MCMC algorithm described in Section~4.4.3, run for 50{,}000 iterations with a burn-in of 10{,}000 iterations and thinning by a factor of 10, yielding 4{,}000 saved posterior samples. The Metropolis--Hastings proposals were tuned so that the acceptance rates for the $\alpha_i$, $\beta_j$, and $\omega_{jl}^{(d)}$ fell within the commonly recommended range of approximately 0.2 to 0.4; the prior hyperparameters, proposal specifications, and full sampler settings are reported in Section~C of the Supplementary Material.


\textcolor{blue}{Convergence was assessed by running multiple independent chains initialized from dispersed starting values, supplemented by within-chain trace plots and effective sample siace (ESS) calculations. All Gelman--Rubin statistics were consistent with convergence, showing no evidence of between-chain disagreement (Table rhat) of the Supplementary Material.The post-burn-in trace plots for the ability ($\alpha_i$) and difficulty ($\beta_j$) parameters exhibited rapid mixing and stable oscillation around posterior mode with no evidence of systematic drift, prolonged transient behavior, or multimodality; because the action-weight parameters are numerous (2{,}025 in total across all items), their convergence was examined both through global summaries and through focused inspection of those weights that enter the reported important-action results, which likewise showed stable sampling without apparent trends. Full trace plots, effective sample sizes, and Gelman--Rubin statistics for all parameter blocks are provided in the Supplementary Material (Table rhat, Table ess, and the accompanying trace-plot figures). On the basis of these diagnostics, we concluded that the posterior sampler had reached its stationary distribution and that the saved samples provided a reliable basis for the substantive interpretation reported in the following subsections.}

% Convergence was assessed using multiple independent chains initialized from dispersed starting values. All Gelman--Rubin statistics were consistent with convergence, showing no evidence of between-chain disagreement. \textcolor{red}{IH: Need to rewrite!}

% Figure(Trace plot) displays representative trace plots for the ability ($\alpha_i$) and difficulty ($\beta_j$) parameters. The sampled values exhibited rapid mixing and stable oscillation around the posterior mode throughout the post-burn-in period, with no evidence of systematic drift, prolonged transient behavior, or multimodality. Because the action-weight parameters are numerous (2{,}025 in total across all items), their convergence was examined both through global summaries and through focused inspection of those weights that enter the reported important-action results; these trace plots showed stable sampling without apparent trends. Full trace plots, effective sample sizes, and Gelman--Rubin statistics for all parameter blocks are provided in Section~C of the Supplementary Material. \textcolor{red}{IH: Need to rewrite! Figure xxx of the Supplementary Material이 되겠지. 여기에는 굳이 trace plot을 넣을 필요 없다. }

%On the basis of these diagnostics, we concluded that the posterior sampler had reached its stationary distribution and that the saved samples provided a reliable basis for the substantive interpretation reported in the following subsections. \textcolor{red}{IH: Combine all three paragraph!}


\subsection{Summary of Selected Actions}

Across all 14 PSTRE items, a total of 126 actions were identified as important under the criterion that the 95\% highest posterior density (HPD) interval of $\Delta E_{jl}$ excludes zero. Table~\ref{tab:selected_actions} reports the number of unique action types in each item's model, the number identified as important, and the directional breakdown of the posterior action effect.

\begin{table}[htb]
\centering
\caption{Summary of important actions identified per item. ``Total'' is the number of unique action types in the model for that item; ``Selected'' is the number identified as important at the 95\% HPD criterion; ``Pos.'' and ``Neg.'' count the selected actions with positive and negative $\Delta E_{jl}$, respectively.}
\label{tab:selected_actions}
\begin{tabular}{llrrrr}
\toprule
Item & Task & Total & Selected & Pos. & Neg. \\
\midrule
PS1-1 & Party Invitations (Part~1) & 119 & 16 & 13 & 3 \\
PS1-2 & Party Invitations (Part~2) & 141 & 10 &  5 & 5 \\
PS1-3 & CD Tally                   & 124 &  3 &  2 & 1 \\
PS1-4 & Sprained Ankle (Part~1)    &  47 &  7 &  6 & 1 \\
PS1-5 & Sprained Ankle (Part~2)    &  48 &  4 &  1 & 3 \\
PS1-6 & Tickets                    & 182 & 16 &  8 & 8 \\
PS1-7 & Class Attendance           & 200 & 11 &  5 & 6 \\
\midrule
PS2-1 & Club Membership (Part~1)   & 156 & 10 &  5 & 5 \\
PS2-2 & Club Membership (Part~2)   & 289 &  4 &  1 & 3 \\
PS2-3 & Book Order                 &  70 &  2 &  2 & 0 \\
PS2-4 & Meeting Room               & 206 & 14 &  6 & 8 \\
PS2-5 & Reply All                  & 112 & 10 &  3 & 7 \\
PS2-6 & Locate Email               & 146 & 11 &  6 & 5 \\
PS2-7 & Lamp Return                & 185 &  8 &  4 & 4 \\
\midrule
Total &                            & 2025 & 126 & 67 & 59 \\
\bottomrule
\end{tabular}
\end{table}

The 126 selected actions represent approximately 6.2\% of the 2{,}025 unique action--item combinations considered in the analysis. The number of selected actions per item ranged from 2 (PS2-3, Book Order) to 16 (PS1-1, Party Invitations Part~1; PS1-6, Tickets), with a median of 10. Of the 126 selected actions, 67 (53.2\%) had positive posterior mean $\Delta E_{jl}$ and 59 (46.8\%) had negative posterior mean. This overall proportion is consistent with the sparse selection structure induced by the spike-and-slab prior. Figure(Delta Panel) shows the posterior distributions of $\Delta E_{jl}$ for the selected actions in each item, arranged by block.


% \begin{figure}[htbp]
% \centering
% % Placeholder: 2x7 panel grid of Delta E_{jl} posteriors per item.
% % Source panels: lstm_ae_deltaE_ps{1,2}_[1-7]_D1.pdf
% \includegraphics[width=\textwidth]{figs/deltaE_panels_D1.pdf}
% \caption{Posterior distributions of $\Delta E_{jl}$ for the selected important actions in each of the 14 PSTRE items, arranged in a $2\times 7$ grid by block. For each item, only actions whose 95\% HPD interval of $\Delta E_{jl}$ excludes zero are displayed.}
% \label{fig:deltaE_panels}
% \end{figure}

\subsection{Item-Level Analysis}

There are several cross-item patterns that emerge from the selected actions and that inform the interpretation framework applied below. First, the directional balance between positive and negative selected actions varies across items in a manner that reflects the functional structure of the task. Several items show a predominantly positive pattern --- for example, PS1-4 (Sprained Ankle Part~1; 6 positive, 1 negative) and PS2-3 (Book Order; 2 positive, 0 negative) --- indicating that correct problem-solving in those tasks is characterized by a distinct set of productive behaviors. Other items show a negative-dominant or balanced pattern --- such as PS2-5 (Reply All; 3 positive, 7 negative) and PS2-4 (Meeting Room; 6 positive, 8 negative) --- indicating that behaviors more prevalent among incorrect respondents also help to distinguish the two groups.

Second, the number of selected actions does not increase monotonically with the size of the action vocabulary. PS2-2 (Club Membership Part~2) has the largest vocabulary (289 unique actions) yet yields only 4 selected actions, whereas PS1-1 and PS1-6 have smaller vocabularies (119 and 182, respectively), each producing 16. This pattern indicates that the sparsity of behaviorally important actions depends on the task structure and the degree to which the interface generates discriminating interactions, rather than on the vocabulary size alone.

Third, we observe a systematic discordance between marginal action frequencies and the sign of $\Delta E_{jl}$ across multiple items. Actions performed predominantly by correct respondents can carry negative partial effects, and actions performed more often by incorrect respondents can carry positive partial effects. This pattern arises because the marginal frequency advantage of an action for correct respondents is largely absorbed by the respondent ability parameter $\alpha_i$. The residual partial effect captured by $\omega_{jl}\bar{C}_{ijl}$ reflects within-ability-group differences in the sequential context in which the action was performed, rather than the total association between the action and response accuracy. This distinction is central to interpreting the case studies below; $\Delta E_{jl}$ should be read as a conditional process-level effect, not as a total effect of performing the action.

Fourth, the interpretive scope of $\Delta E_{jl}$ depends on the richness of the problem-solving sequence. In items with extended, multi-step action sequences, the LSTM encoder captures contextual information --- such as exploration order, trial-and-error patterns, and temporal pacing --- that is partly independent of respondent ability. In these items, the sign and magnitude of $\Delta E_{jl}$ can support substantive interpretation of behavioral differences between correct and incorrect respondents. In items with short, simple sequences where the only meaningful behavioral variation is the identity of the final choice, the latent representation $\bar{C}_{ijl}$ co-varies strongly with $\alpha_i$, and the partial effect may not carry a stable process-level interpretation. For this reason, we distinguish three item types based on where behavioral discrimination is concentrated, and we calibrate the depth of interpretation accordingly.

Based on the functional character of their selected actions, the 14 PSTRE items can be organized into three types. In \textit{outcome-dominant} items, most selected actions correspond to final-choice operations, and the partial effects are interpreted primarily at the level of selection significance rather than process-level behavioral differences. In \textit{process-dominant} items, selected actions are distributed across intermediate behavioral steps, supporting richer interpretation of how correct and incorrect respondents differ throughout the solution process. In \textit{hybrid} items, both outcome-proximal and process-level actions appear among the selected set. Table~\ref{tab:item_types} summarizes the classification.

\begin{table}[htb]
\centering
\caption{Classification of PSTRE items by the functional distribution of selected important actions.}
\label{tab:item_types}
\hspace*{-0.9cm}
\begin{tabular}{cll}
\toprule
Label & Description & Items \\
\midrule
Outcome-dominant & Most selected actions are final-choice operations & PS1-3, PS1-4, PS1-5, PS1-7, PS2-3 \\
Process-dominant & Selected actions span intermediate process steps & PS1-1, PS1-2, PS2-1, PS2-5, PS2-6 \\
Hybrid           & Outcome-proximal and process actions coexist     & PS1-6, PS2-2, PS2-4, PS2-7 \\
\bottomrule
\end{tabular}
\end{table}

The following subsections present one representative case study per type: PS1-3 (CD Tally) for outcome-dominant items, PS1-1 (Party Invitations Part~1) for process-dominant items, and PS2-7 (Lamp Return) for hybrid items. For outcome-dominant items, interpretation is restricted to the first-order level --- the selection of the action itself as carrying information beyond ability and difficulty --- without extending to claims about the direction of the partial effect, because the simple sequence structure in these items limits the process-level information encoded in $\bar{C}_{ijl}$. For process-dominant and hybrid items, both first-order and second-order interpretations are applied, including the relationship between marginal frequency patterns and the sign of $\Delta E_{jl}$. Detailed results for the remaining items are provided in Section~D of the Supplementary Material.

\subsubsection{Outcome-Dominant Items: CD Tally (PS1-3)}

The CD Tally task asks respondents to count the number of Blues-genre compact discs listed in a sortable spreadsheet and enter the total using a numeric dropdown menu. This task is representative of outcome-dominant items in that meaningful behavioral variation is concentrated almost entirely at the final selection step, with limited discriminative information in intermediate navigation actions.

\begin{table}[htb]
\centering
\caption{Important actions for PS1-3 (CD Tally). $n_C$ and $n_I$ denote the number of respondents in the correct and incorrect groups who performed the action at least once.}
\label{tab:ps13_actions}
\small
\begin{tabular}{clrrl}
\toprule
Dir. & Action & $n_C$ & $n_I$ & $\Delta E$ (95\% HPD) \\
\midrule
$(+)$ & \texttt{combobox\_menulist\_index9}     & 490 &   2 & $\phantom{-}$0.2726 (0.2457, 0.3009) \\
$(+)$ & \texttt{combobox\_sortablecol1\_index3} & 199 &  20 & $\phantom{-}$0.0005 (0.0000, 0.0010) \\
$(-)$ & \texttt{toolbar\_webapp}               & 479 & 196 & $-$0.0037 ($-$0.0064, $-$0.0008) \\
\bottomrule
\end{tabular}
\end{table}

Three actions were identified as important for this item; the correct answer entry, a sorting strategy, and an interface switch. At the first-order level, all three carry information about response accuracy beyond what is explained by respondent ability and item difficulty. The action \texttt{combobox\_menulist\_index9} corresponds to selecting 8 as the Blues CD count, which is the correct answer. It was performed by 490 correct respondents and only 2 incorrect respondents. This extreme frequency imbalance constitutes quasi-complete separation in the logistic component of the model, and the resulting $\Delta E = 0.273$ is likely inflated in magnitude. The direction of the effect and its clear separation from zero are interpretable, but the magnitude should not be compared directly with those of other actions.

The action \texttt{combobox\_sortablecol1\_index3} corresponds to sorting the spreadsheet by the genre column, a strategy that facilitates counting by grouping all Blues titles together. This action reflects a deliberate problem-solving strategy rather than a final answer selection: respondents who sorted the spreadsheet adopted a more efficient approach to the counting task. Its positive $\Delta E = 0.001$, though small, suggests that this sorting strategy is associated with a modestly favorable action context after controlling for respondent ability. Among the three selected actions for this item, this is the only one that admits a partial second-order interpretation, as it represents a procedural choice rather than an outcome-defining operation.

The action \texttt{toolbar\_webapp} represents switching from the spreadsheet interface to the web application, an operation that is not required to complete the task. Although this action was performed more frequently by correct respondents ($n_C = 479$) than by incorrect respondents ($n_I = 196$), its $\Delta E = -0.004$ is negative. This discordance between the marginal frequency pattern and the partial effect illustrates the role of the ability parameter in absorbing the total association: the positive marginal frequency advantage for correct respondents is captured by $\alpha_i$, and the residual partial effect reflects the fact that, within ability-matched groups, unnecessary interface switching is associated with a less favorable action context.

This item illustrates both the strengths and the interpretive boundaries of the proposed framework when applied to outcome-dominant tasks. The model identifies the correct answer selection without requiring explicit labeling of which interface element corresponds to the correct response, and it detects a strategic sorting behavior that contributes process-level information. At the same time, the predominance of final-choice actions among the selected set, together with the quasi-complete separation in the correct answer action, limits the extent to which the partial effects can support detailed process-level interpretation.

% Supp의 해석으로 넘기기
% Another outcome-dominant item, PS1-5 (Sprained Ankle Part~2), exhibits a more extreme version of this limitation. All four selected actions in PS1-5 are combobox choices corresponding to specific answer options, and no intermediate process action was identified as important. In this item, the correct answer action (\texttt{combobox\_index2}; $n_C = 661$, $n_I = 10$) carries a negative $\Delta E = -0.035$ despite its overwhelming frequency advantage for correct respondents. This systematic discordance between marginal frequency and partial effect across all selected actions arises because the simple procedural structure of the task---respondents view web pages and then select an answer---limits the latent representation to the identity of the chosen response, which co-varies strongly with respondent ability. For items with this structure, interpretation is accordingly restricted to the first-order level: the selected actions carry information beyond ability and difficulty, but the sign of $\Delta E$ does not support stable inference about behavioral strategies.

\subsubsection{Process-Dominant Items: Party Invitations (PS1-1)}

The Party Invitations task requires respondents to read incoming email messages and sort each sender into either a ``Can Come'' or ``Cannot Come'' folder based on the stated availability. The task is scored polytomously (0--3), with the binary outcome defined as a score of 2 or 3. Sixteen actions were identified as important, the largest count among all 14 items. This item is representative of process-dominant items in that the selected actions span the full breadth of the problem-solving process, from message reading and folder manipulation to non-task-relevant exploration. Both first-order and second-order interpretations are applicable, as the action sequences for this item are sufficiently rich (median sequence length 16; IQR 11--25) to provide latent representations that encode process-level context beyond the identity of individual actions.

\begin{table}[htb]
\centering
\caption{Selected important actions for PS1-1 (Party Invitations Part~1). Eight representative actions are shown; the complete list of all 16 selected actions is provided in Section~D of the Supplementary Material. $n_C$ and $n_I$ as in Table~\ref{tab:ps13_actions}.}
\label{tab:ps11_actions}
\small
\begin{tabular}{clrrl}
\toprule
Dir. & Action & $n_C$ & $n_I$ & $\Delta E$ (95\% HPD) \\
\midrule
$(+)$ & \texttt{mail-drop\_u01a-cancomefolder}     & 783 & 130 & $\phantom{-}$0.3263 ($\phantom{-}$0.2121, $\phantom{-}$0.4439) \\
$(+)$ & \texttt{mail-drag\_u01a-item104}           & 772 &  58 & $\phantom{-}$0.1372 ($\phantom{-}$0.0759, $\phantom{-}$0.1956) \\
$(+)$ & \texttt{mail-drop\_u01a-cannotcomefolder}  & 662 &  33 & $\phantom{-}$0.1029 ($\phantom{-}$0.0494, $\phantom{-}$0.1554) \\
$(+)$ & \texttt{mail-drop\_u01a-partyfolder}       &  58 &  70 & $\phantom{-}$0.1044 ($\phantom{-}$0.0351, $\phantom{-}$0.1771) \\
$(+)$ & \texttt{mail-drag\_u01a-item201}           &  34 &  88 & $\phantom{-}$0.0273 ($\phantom{-}$0.0104, $\phantom{-}$0.0445) \\
$(+)$ & \texttt{folder-viewed\_u01a-cannotcomefolder} & 346 & 188 & $\phantom{-}$0.0000 ($\phantom{-}$0.0000, $\phantom{-}$0.0000) \\
$(-)$ & \texttt{folder-viewed\_u01a-inboxfolder}  & 334 & 150 & $-$0.0232 ($-$0.0351, $-$0.0110) \\
$(-)$ & \texttt{toolbar\_mailapp}                 &  17 &  54 & $-$0.0004 ($-$0.0007, $-$0.0001) \\
\bottomrule
\end{tabular}
\end{table}

At the first-order level, the 16 selected actions encompass the core operations required for the task---reading target messages (e.g., \texttt{mail-viewed\_item105}, \texttt{mail-viewed\_item102}), dragging messages for classification (e.g., \texttt{mail-drag\_item104}, \texttt{mail-drag\_item101}), and dropping messages into the designated folders (\texttt{mail-drop\_cancomefolder}, \texttt{mail-drop\_cannotcomefolder})---as well as non-task-relevant exploration behaviors such as viewing extraneous messages (\texttt{mail-viewed\_item202}), opening the parent folder (\texttt{folder-viewed\_partyfolder}), and switching interfaces (\texttt{toolbar\_mailapp}). The selection of both task-relevant and non-task-relevant actions provides evidence that the model captures process-level information extending beyond the identity of the correct classification operations.

The largest positive effects are associated with the classification steps that directly correspond to the task requirements: dropping messages into the Can Come folder ($\Delta E = 0.326$), dragging a task-relevant message ($\Delta E = 0.137$), and dropping messages into the Cannot Come folder ($\Delta E = 0.103$). For these actions, the marginal frequency pattern and the partial effect are concordant---correct respondents performed them more frequently and the partial effect is positive---indicating that the process context in which these actions were performed provides additional discriminative information beyond the contribution already absorbed by respondent ability.

Several selected actions exhibit a discordance between their marginal frequency pattern and the sign of $\Delta E_{jl}$, which can be interpreted within the partial-effect framework described above. The actions \texttt{mail-drag\_u01a-item201} ($n_C = 34$, $n_I = 88$; $\Delta E = +0.027$) and \texttt{mail-drag\_u01a-item202} ($n_C = 36$, $n_I = 84$; $\Delta E = +0.002$) involve messages that are not classification targets for this task, and both were performed more often by incorrect respondents. The negative marginal frequency association is absorbed by the ability parameter. The positive residual partial effect may reflect the fact that, among respondents of comparable ability, those who engaged with these extraneous messages as part of a systematic exploration pattern---examining all available messages before making classification decisions---experienced a more favorable action context than those who engaged with the same messages in a less structured manner. This interpretation is consistent with the sequential context encoded by the LSTM, but should be regarded as tentative given the relatively small number of correct respondents who performed these actions.

The action \texttt{folder-viewed\_u01a-inboxfolder} ($n_C = 334$, $n_I = 150$; $\Delta E = -0.023$) exhibits the reverse pattern: it was performed more often by correct respondents, yet its partial effect is negative. Opening the inbox folder is the default starting point for the task and is performed by a large proportion of all respondents. The positive marginal frequency association with correct response is absorbed by the ability parameter. The negative residual partial effect suggests that, among respondents of comparable ability, repeated returns to the inbox folder may reflect less efficient navigation behavior, such as re-checking messages that have already been read rather than proceeding with classification.

The action \texttt{mail-drop\_u01a-partyfolder} ($n_C = 58$, $n_I = 70$; $\Delta E = +0.104$) represents placing a message into the parent Party folder rather than the required Can Come or Cannot Come subfolders. Nearly equal numbers of correct and incorrect respondents performed it. Its positive partial effect suggests that, among respondents of comparable ability, those who placed messages in the parent folder did so in a sequential context associated with eventual task success---for example, using the parent folder as a temporary holding location during the classification process before redistributing messages to the appropriate subfolders. This interpretation is plausible but not directly verifiable from the model output alone.

Overall, this item illustrates the richer interpretive potential of process-dominant tasks. The extended action sequences allow the LSTM encoder to capture contextual information that is partially independent of respondent ability, and the resulting partial effects reveal behavioral distinctions---such as the difference between systematic and unsystematic exploration of non-target messages---that would not be visible from action frequencies or response outcomes alone.

\subsubsection{Hybrid Items: Lamp Return (PS2-7)}

The Lamp Return task requires respondents to initiate a return of a misdirected lamp by navigating between a simulated email client and a web form. The correct procedure involves requesting a return authorization number via an email link, retrieving the number from the subsequent reply, and submitting the completed return request through the web interface. Eight actions were identified as important, spanning both procedural intermediate steps and final confirmatory operations. The action sequences for this item are moderately long (median sequence length 15; IQR 6--24), providing sufficient sequential context for both first-order and second-order interpretation.

\begin{table}[htb]
\centering
\caption{Important actions for PS2-7 (Lamp Return). $n_C$ and $n_I$ as in Table~\ref{tab:ps13_actions}. Actions are listed in approximate procedural order.}
\label{tab:ps27_actions}
\small
\begin{tabular}{clrrl}
\toprule
Dir. & Action & $n_C$ & $n_I$ & $\Delta E$ (95\% HPD) \\
\midrule
$(-)$ & \texttt{textlink\_u023-pg2\_txt21\_self}    & 252 &  13 & $-$0.0032 ($-$0.0046, $-$0.0018) \\
$(-)$ & \texttt{combobox\_u023-pg2\_menu1\_index1}  & 541 &  18 & $-$0.0171 ($-$0.0329, $-$0.0005) \\
$(-)$ & \texttt{textbox-onfocus\_u023-pg2\_txt19}  & 519 &   9 & $-$0.0340 ($-$0.0472, $-$0.0224) \\
$(+)$ & \texttt{mail-viewed\_u23-item305}          & 480 & 132 & $\phantom{-}$0.0010 ($\phantom{-}$0.0003, $\phantom{-}$0.0016) \\
$(+)$ & \texttt{toolbar\_back-btn}                & 502 & 233 & $\phantom{-}$0.0353 ($\phantom{-}$0.0166, $\phantom{-}$0.0545) \\
$(+)$ & \texttt{button\_u023-pg2\_txt22}           & 512 &  11 & $\phantom{-}$0.0505 ($\phantom{-}$0.0397, $\phantom{-}$0.0615) \\
$(+)$ & \texttt{keypress2}                        & 345 &   9 & $\phantom{-}$0.4619 ($\phantom{-}$0.1941, $\phantom{-}$0.7227) \\
$(-)$ & \texttt{textlink\_u023-default\_txt5}      & 150 & 324 & $-$0.0030 ($-$0.0055, $-$0.0005) \\
\bottomrule
\end{tabular}
\end{table}

At the first-order level, the eight selected actions cover each major stage of the multi-step return procedure: requesting the authorization email (\texttt{textlink\_txt21\_self}), selecting the return reason (\texttt{combobox\_menu1\_index1}), activating the input field for the authorization number (\texttt{textbox-onfocus\_txt19}), reviewing the authorization email (\texttt{mail-viewed\_item305}), navigating back to the return form (\texttt{toolbar\_back-btn}), entering the authorization number (\texttt{keypress2}), and submitting the final request (\texttt{button\_txt22}). The breadth of this selection indicates that the model identifies discriminative information at multiple points along the procedural sequence, rather than concentrating it at a single outcome-defining step. This is the defining characteristic of hybrid items: both intermediate process actions and final confirmatory operations contribute information beyond respondent ability and item difficulty.

At the second-order level, a systematic sign pattern emerges when the selected actions are arranged in approximate procedural order. The three actions associated with the earlier stages of the procedure---requesting the authorization email (\texttt{textlink\_txt21\_self}; $n_C = 252$, $n_I = 13$), selecting the return reason (\texttt{combobox\_menu1\_index1}; $n_C = 541$, $n_I = 18$), and activating the authorization number input field (\texttt{textbox-onfocus\_txt19}; $n_C = 519$, $n_I = 9$)---all carry negative $\Delta E$ despite being performed predominantly by correct respondents. This discordance between marginal frequency and partial effect is consistent with the absorption mechanism: the positive total association between these early-stage actions and correct response is captured by the ability parameter, and the residual partial effect reflects within-ability-group differences in the sequential context surrounding these actions.

By contrast, the four actions associated with the later stages of the procedure---reviewing the authorization email (\texttt{mail-viewed\_item305}; $\Delta E = +0.001$), navigating back to the return form (\texttt{toolbar\_back-btn}; $\Delta E = +0.035$), entering the authorization number (\texttt{keypress2}; $\Delta E = +0.462$), and submitting the final request (\texttt{button\_txt22}; $\Delta E = +0.051$)---all carry positive partial effects.

This anterior--posterior sign differentiation may be interpreted as follows. The early procedural steps are initiated by a large proportion of respondents, including many who attempt the return procedure but do not complete it. Among respondents of comparable ability, reaching the early stages does not by itself differentiate those who ultimately succeed from those who do not; accordingly, the residual partial effect at these stages does not favor the correct-response group. In contrast, the later stages---retrieving the authorization number from the reply email, navigating back to the form, entering the number, and submitting the request---are reached primarily by respondents who sustain the multi-step procedure to completion. The positive partial effects at these stages suggest that, within ability-matched groups, the process context of completing the later procedural steps carries an incremental association with task success that is not fully explained by respondent ability. This pattern is consistent with the interpretation that correct and incorrect respondents diverge not at the point of initiating the procedure but at the point of completing it.

Two individual actions warrant additional comment. The action \texttt{keypress2}, corresponding to typing the return authorization number, carries the largest $\Delta E$ among the eight selected actions ($\Delta E = 0.462$; 95\% HPD: 0.194, 0.723). However, the extreme frequency imbalance ($n_C = 345$, $n_I = 9$) constitutes quasi-complete separation, and the wide HPD interval reflects the resulting estimation uncertainty. The direction of the effect is interpretable, but the magnitude should not be compared with those of other actions in this item. The action \texttt{textlink\_u023-default\_txt5} ($n_C = 150$, $n_I = 324$; $\Delta E = -0.003$) was performed predominantly by incorrect respondents and carries a negative partial effect; however, the specific interface function of this link could not be conclusively identified from the available log metadata, so we refrain from offering a substantive behavioral interpretation for this action.

This item illustrates the analytical value of the hybrid classification. The coexistence of outcome-proximal actions (authorization number entry, final submission) and intermediate process actions (email request, reason selection, form navigation) within the same selected set allows the model to reveal where in a multi-step procedure the behavioral trajectories of correct and incorrect respondents diverge. The systematic anterior--posterior sign pattern provides a more informative characterization of the problem-solving process than would be available from either final-step analysis or action frequency comparison alone.

%%%%%%%%%%%%%%%%%%%%%%%%%%%%%%%%%%%%%%%%%%%%%%%%
\section{Simulation Studiess}

\subsection{Experiment Settings}

We conducted a simulation study to evaluate the variable-selection performance of the proposed representation-learning action-IRT model with a spike-and-slab prior, specifically whether the model recovers a known set of important actions under data-generating conditions that approximate the empirical application. Because the framework treats the LSTM autoencoder output as a fixed covariate, the simulation targets the final representation-learning action-IRT stage: the observed respondent--item structure, action-occurrence patterns, and reduced latent action values are held fixed at their empirical values, and only the binary response outcomes are regenerated. This parametric-bootstrap design ensures that the simulation inherits the covariate structure, missingness pattern, and sparsity level of the actual analysis.

The data-generating model is calibrated to the empirical representation-learning action-IRT fit reported in Section~5: the model is first fitted to the PIAAC data, and the resulting posterior then defines the ground truth for the simulation. The simulated data preserve the structure of the empirical analysis---14 PSTRE items, 1{,}996 respondents, 2{,}025 item-specific action types, and a one-dimensional latent representation ($D=1$)---with the covariate values $\bar{C}_{ijl}$ held at their empirical values. The number of unique actions per item and the corresponding counts of true important actions are listed in Section~C of the Supplementary Material.

True parameter values were taken from the empirical posterior: respondent ability and item difficulty were fixed at their posterior means, $\alpha_i^{\mathrm{true}}=\mathbb{E}(\alpha_i\mid\mathbf{Y})$ and $\beta_j^{\mathrm{true}}=\mathbb{E}(\beta_j\mid\mathbf{Y})$. The 126 action--item combinations identified as important by the empirical fit (6.22\% of all 2{,}025 combinations; Section~5) were designated as the true important actions, each with $\lambda_{jl}^{\mathrm{true}}=1$ and true weight equal to the posterior mean of $\omega_{jl}$ under slab membership. The remaining actions were assigned $\lambda_{jl}^{\mathrm{true}}=0$ with weights drawn from the spike component; the exact draw is specified in Section~C of the Supplementary Material.

For each replication, binary responses were generated for the observed respondent--item pairs as $Y_{ij}\sim\mathrm{Bernoulli}\bigl(\mathrm{logistic}(\eta_{ij})\bigr)$, with
\[
\eta_{ij}
=
\alpha_i^{\mathrm{true}}+\beta_j^{\mathrm{true}}
+\sum_{l\in A_{ij}}\omega_{jl}^{\mathrm{true}}\,\bar{C}_{ijl},
\]
where $A_{ij}$ is the set of unique actions performed by respondent $i$ on item $j$. We carried out five independent replications with distinct random seeds, applying the same MCMC sampler and posterior-sampling settings used in the empirical analysis to each simulated response matrix (Section~5.1; full settings in Section~C of the Supplementary Material).

Variable-selection performance was evaluated using the posterior inclusion probability, $\mathrm{PIP}_{jl}=P(\lambda_{jl}=1\mid\mathbf{Y})$, estimated as the proportion of saved iterations in which $\lambda_{jl}=1$; because $D=1$, the action-level PIP coincides with the single-dimension PIP. We report the area under the receiver operating characteristic curve (AUC), computed with the true inclusion indicator as the label and PIP as the score, and the sensitivity under the threshold rule $\mathrm{PIP}_{jl}\geq 0.5$.

\subsection{Results}

Table~\ref{tab:simulation_recovery} summarizes variable-selection performance across the five replications. The simulated correct-response rate was stable (mean 0.4300, SD 0.0028), indicating that the data-generating mechanism produced consistent response distributions across replications. The action-level AUC averaged 0.9431 (range 0.9257--0.9550), and sensitivity under the $\mathrm{PIP}\geq 0.5$ threshold averaged 0.9349 (range 0.9127--0.9444). All five replications yielded AUC above 0.92 and sensitivity above 0.91, indicating that the result was not driven by a single favorable run.

\begin{table}[htbp]
\centering
\caption{Simulation performance across five replications.}
\label{tab:simulation_recovery}
\begin{tabular}{lS[table-format=1.4]S[table-format=1.4]S[table-format=1.4]S[table-format=1.4]S[table-format=1.4]}
\toprule
Metric & {Mean} & {SD} & {Median} & {Min.} & {Max.} \\
\midrule
Simulated correct rate            & 0.4300 & 0.0028 & 0.4303 & 0.4269 & 0.4329 \\
Action AUC                        & 0.9431 & 0.0108 & 0.9460 & 0.9257 & 0.9550 \\
Action sensitivity, PIP $\geq 0.5$ & 0.9349 & 0.0130 & 0.9365 & 0.9127 & 0.9444 \\
\bottomrule
\end{tabular}
\end{table}

The high AUC values indicate that the PIP reliably separates truly important actions from non-important ones, and the high sensitivity confirms that the model recovers the vast majority of true important actions across replications. These results provide evidence that the spike-and-slab representation-learning action-IRT model can recover the empirically grounded set of important actions with high discrimination when the observed action structure and LSTM latent values are preserved.

\section{Conclusion}

This paper addressed the problem of identifying behaviorally important actions in problem-solving process data from large-scale computer-based assessments, where existing approaches largely rely on analyst-selected process indicators that depend on prior knowledge of the task structure. We proposed an end-to-end, three-stage framework that (i)~embeds each action unit through a hybrid Word2Vec model combining skip-gram training over action sequences with FastText-style composition over action tokens, (ii)~reduces each respondent's temporally augmented action sequence to a low-dimensional latent representation through an LSTM autoencoder that preserves sequential ordering, and (iii)~incorporates the resulting latent action values into an extended Rasch (representation-learning action-IRT) model with a spike-and-slab prior. The prior performs automatic, data-driven selection of actions associated with response accuracy beyond respondent ability and item difficulty, without pre-specifying which actions are expected to be discriminative.

Applied to log process data from 14 PSTRE items in the United States PIAAC study, the framework selected 126 important actions among 2,025 unique action--item combinations (a selection rate of approximately 6.2\%), spanning the full range of observed problem-solving behavior. Item-level case studies revealed behavioral patterns that are not visible from response outcomes or raw action frequencies alone and motivated a typological distinction among outcome-dominant, process-dominant, and hybrid items: behavioral discrimination is concentrated at the final choice step in outcome-dominant items (e.g., CD Tally, PS1-3), distributed across intermediate steps in process-dominant items (e.g., Party Invitations, PS1-1), and organized into a systematic anterior--posterior sign pattern in hybrid items (e.g., Lamp Return, PS2-7) that locates where correct and incorrect respondents diverge. Because high-frequency actions are largely absorbed by the ability parameter, these ability-adjusted partial effects identify discriminative behaviors that raw frequencies would misrepresent.

Several limitations qualify these results. The process-level interpretation of the action effects is weaker for structurally simple, outcome-dominant items, where the latent action representation co-varies strongly with respondent ability; actions performed by only one response group are subject to quasi-complete separation, so their effect magnitudes are not directly comparable; the analysis uses a single latent dimension ($D=1$) and fits the embedding and autoencoder separately for each item, which limits cross-item information sharing.

Notwithstanding these limitations, the proposed framework offers a principled and reproducible procedure for extracting behavioral insight from process data within a well-established psychometric model, yielding identified important actions that can serve as a data-driven complement to expert-designed process indicators for post-hoc strategy assessment. Promising directions for future work include sensitivity analyses across latent dimensions and prior specifications, comparison with alternative action representations, extension to multi-dimensional and cross-item joint modeling, and integration with response time and other behavioral signals toward a richer account of problem-solving competence from digital assessment data.

\section*{Acknowledgments}

This work was partially supported by the National Research Foundation of Korea [grant number NRF-2021S1A3A2A03088949, RS-2023-00217705, and RS-2024-00333701; Basic Science Research Program awarded to I.H.J.] and the ICAN (ICT Challenge and Advanced Network of HRD) support program [grant number RS-2023-00259934], supervised by the IITP (Institute of Information \& Communications Technology Planning \& Evaluation), and the Ministry of Trade, Industry, and Energy (MOTIE), Korea, under the project ``Industrial Technology Infrastructure Program'' (RS-2024-00466693). Correspondence should be addressed to Ick Hoon Jin, Department of Applied Statistics, Department of Statistics and Data Science, Yonsei University, Seoul, Republic of Korea. E-Mail: ijin@yonsei.ac.kr. 

\newpage
%\bibliographystyle{Chicago}
\bibliographystyle{apalike}
\bibliography{reference}

\end{document}